\subsection{由伪度量族定义的一致结构的性质}

设 \(\mathcal{U}\) 是在集合 \(X\) 上由一族有限伪度量 \((f_{i})\) 定义的一致结构。若给 \(X \times X\) 赋以 \(\mathcal{U}\) 与其自身的乘积一致结构，则每个实值函数 \(f_{i}\) 在 \(X \times X\) 上都一致连续；因为由 (1) 我们有

\[
\left| f _ {t} (x, y) - f _ {t} \left(x ^ {\prime}, y ^ {\prime}\right) \right| \leqslant f _ {t} (x, x ^ {\prime}) + f _ {t} (y, y ^ {\prime}),
\]

因而关系 \(f_{i}(x, x') \leqslant \varepsilon / 2\) 与 \(f_{i}(y, y') \leqslant \varepsilon / 2\) 蕴含

\[
\left| f _ {t} (x, y) - f _ {t} \left(x ^ {\prime}, y ^ {\prime}\right) \right| \leqslant \varepsilon .
\]

\(\mathcal{U}\) 为豪斯多夫的必要充分条件（由 \(\mathcal{U}\) 的围域的定义）是：对每对不同的点 \(x, y\) （属于 \(\mathbf{X}\) ），存在指标 \(\iota\) 使得 \(f_{\iota}(x, y) \neq 0\) 。

特别地，若 \(\mathcal{U}\) 由单个伪度量 \(f\) 定义，则 \(\mathcal{U}\) 是豪斯多夫的当且仅当关系 \(f(x,y)=0\) 蕴含 \(x=y\) （参看§ 2）。若 \(\mathcal{U}\) 不是豪斯多夫的，则 \(\mathcal{U}\) 的所有围域之交是 \(\mathbf{X} \times \mathbf{X}\) 的一个子集，它由点对 \((x,y)\) （满足 \(f_{i}(x,y)=0\) 对一切 \(\iota\) ）组成；这个子集是等价关系 \(\mathbf{R}\) 在 \(\mathbf{X}\) 上的图象，而与 \(\mathcal{U}\) 相伴的豪斯多夫一致结构定义在 \(\mathbf{X}/\mathbf{R}\) 上（参看第二章，§ 3，第 8 号）。于是容易验证：函数 \(f_{i}\) （关于 \(x\) 与 \(y\) ）与关系 \(\mathbf{R}\) 相容（《集合论》，\(\mathbf{R}\) ，§ 5，第 7 号），而函数 \(\overline{f_{i}}\) —— 由 \(f_{i}\) 取商（关于 \(x\) 与 \(y\) ）而得 —— 是 \(\mathbf{X}/\mathbf{R}\) 上的伪度量，它们定义与 \(\mathcal{U}\) 相伴的豪斯多夫一致结构（参看§ 2，第 1 号）。

若 A 是 X 的非空子集，则 X 上伪度量在 \(A \times A\) 上的限制显然是 A 上的伪度量。U 在 A 上诱导的一致结构显然就是由到 \(A \times A\) 上的限制（诸伪度量 \(f_{t}\) 的限制）所组成的族定义的一致结构。

现在来考察一致空间 \(\mathbf{X}\) 当 \(\mathcal{U}\) 为豪斯多夫一致结构时的完备化。

命题 1. 设 X 是豪斯多夫一致空间，其一致结构 U 由一族有限伪度量 ( \(f_{t}\) ) 定义，并设 \(\hat{X}\) 是 X 的完备化。则诸函数 \(f_{t}\) 可以用连续性扩张到 \(\hat{\mathbf{X}} \times \hat{\mathbf{X}}\) ；扩张后的函数 \(\overline{f_t}\) 是 \(\hat{\mathbf{X}} \times \hat{\mathbf{X}}\) 上的有限伪度量，且族 \((f_{t})\) 定义 \(\hat{\mathbf{X}}\) 的一致结构。

首先，诸 \(f_{t}\) 可以用连续性扩张到 \(\hat{\mathbf{X}}\times\hat{\mathbf{X}}\) ，因为它们在 \(\mathbf{X}\times\mathbf{X}\) 上一致连续；且扩张后的函数 \(\overline{f_{t}}\) 在 \(\hat{\mathbf{X}}\times\hat{\mathbf{X}}\) 上一致连续（第二章，§ 3，第 6 号，定理 2）；此外，由不等式扩张原理（第四章，§ 5，第 2 号，定理 1），它们是 \(\hat{\mathbf{X}}\) 上的伪度量。用 \(\mathcal{U}_{1}\) 表示由完备化得到的 \(\hat{\mathbf{X}}\) 上的一致结构，用 \(\mathcal{U}_{2}\) 表示由伪度量族 ( \(\overline{f_{t}}\) ) 定义的一致结构。则 \(\mathcal{U}_{2}\) 粗于 \(\mathcal{U}_{1}\) ；因为每个 \(\overline{f_{t}}\) 在 \(\hat{\mathbf{X}}\times\hat{\mathbf{X}}\) 上关于 \(\mathcal{U}_{1}\) 一致连续，故对每个 \(a>0\) 存在 \(\mathcal{U}_{1}\) 的一个围域 V，使得只要 \((x,y)\in\mathrm{V}\) ，就有 \(|\overline{f_{t}}(x,y)-\overline{f_{t}}(x,x)|\leqslant a\) ，即{[}由于 \(\overline{f_{t}}(x,x)=0\) {]}，\(\mathrm{V}\subset\overline{f_{t}}([0,a])\) ；因而 \(\mathcal{U}_{2}\) 的每个围域都是 \(\mathcal{U}_{1}\) 的围域。另一方面，\(\mathcal{U}_{1}\) 与 \(\mathcal{U}_{2}\) 在 X 上诱导同一个一致结构 \(\mathcal{U}\) 。由于 \(\hat{\mathbf{X}}\) 关于 \(\mathcal{U}_{1}\) 完备，由此可知 \(\mathcal{U}_{1}\) 与 \(\mathcal{U}_{2}\) 相同（第二章，§ 3，第 7 号，命题 14）。

